% Source PDF pp. 37--46. The first continued sentence belongs to en-06.
Proceeding as in the proof of 6.2.4, one shows that $U'_1$ is invariant in $E$. After passing to the quotient by $U'_1$, one reduces to the case where $U=\Ga$. The $S$-scheme $E$ (where $S=U$) is then a torsor under the $S$-group $\Gm\times_k S$, hence has a section, since $\Pic(S)=0$ (Exp. VIII 4.3). The extension $E$ is then trivial by A) i).

\par\noindent\emph{Proof of 6.1.1 A) iii) ($U$ smooth, $H$ étale).}
Suppose first that $U$ is connected. The group $U_{\bar k}$ then has a composition series with successive quotients isomorphic to $\Ga$, so $E_{\bar k}$ is trivial by A) ii). Since $H$ is étale, it is clear that the unique lifting of $U_{\bar k}$ in $E_{\bar k}$ is the connected component of $E_{\bar k}$, so this lifting is already defined over $k$. In the general case, after passing to the quotient by the connected component of $E$, one reduces to the case where $E$ is étale, then to the case where $E$ is completely split (6.2.3). Since $U(k)$ is a $p$-group and $H(k)$ has order prime to $p$, one may take as lifting of $U(k)$ the Sylow $p$-subgroup of $E(k)$.

\par\noindent\emph{Proof of 6.1.1 A) iv) ($U$ smooth, $k$ perfect).}
If $U$ is connected, $U$ has a composition series with successive quotients isomorphic to $\Ga$ (4.1.2 b)), and one applies A) ii). In the general case, the foregoing allows us to reduce to the case where $U$ is étale, then to the case where $U$ is completely split and $H$ diagonalizable (6.2.3). Using now a characteristic composition series of $H$ ($H\supset H^0\supset {}_{F^n}(H)\supset0$), one reduces to the case where $H$ is of one of the following three types:

\noindent a) $H$ is étale.

\noindent b) $H=(\Gm)^r$.

\noindent c) $H$ is radicial.

In case a), one applies A) iii); in case c), one applies 1.6. Finally, in case b), one remarks that by Hilbert's Theorem 90, $E\to U$ has a section, so that it suffices to show that $H^2(U,(\Gm)^r)=H^2(U(k),(\Gm)^r(k))=0$. Now $U(k)$ is a finite $p$-group, whereas $(\Gm)^r(k)$ is uniquely $p$-divisible (since $k$ is perfect).

\numpara{6.3}\label{XVII.6.3}\textbf{Proof of 6.1.1 B) and C).}
By 6.2.1, 6.2.3 and 6.2.4, one sees that it suffices to prove B) when $H=\Gm$. One therefore has an exact sequence
\[
1\longrightarrow\Gm\longrightarrow E\longrightarrow\balpha_p\longrightarrow1.
\]
Since $\Gm$ is smooth, one deduces an exact sequence of Lie $p$-algebras (App. II 3.2):
\[
0\longrightarrow\Lie\Gm\longrightarrow\Lie E\longrightarrow\Lie\balpha_p\longrightarrow0.\tag{*}
\]
The group $\balpha_p$ being of height $1$, the extension $E$ is trivial if and only if (App. II 2.2) the exact sequence (*) of Lie $p$-algebras is split. One knows that $\Lie\balpha_p$ is generated by an element $X\ne0$ such that $X^{(p)}=0$ (App. II 2.1); it therefore suffices to show that $X$ lifts to an element $Z$ of $\Lie E$ such that $Z^{(p)}=0$. Now by 5.8.2 ii), there exists a unipotent element $Z$ of $\Lie E$ lifting $X$. Since the unipotent part of $\Lie E$ clearly has dimension at most $1$, one necessarily has $Z^{(p)}=0$.

\par\noindent\emph{Proof of 6.1.1 C).}
If $U'_1$ is a lifting in $E$ of an algebraic subgroup $U_1$ of $U$, $U'_1$ is invariant in $E$, since $E$ is supposed commutative, and we may pass to the quotient by $U'_1$. Using a composition series of $U$ (3.5 ii)), the preceding remark allows us to reduce to the case where $U$ is smooth or equal to $\balpha_p$. But then $E$ is trivial by A) iv) and B).

% Source PDF p. 38.
\numpara{6.4}\label{XVII.6.4}\textbf{Examples of extensions of a unipotent group $U$ by a group of multiplicative type $H$ which are not trivial.}
In view of 6.1.1 A) iv), the problem arises only in characteristic $p>0$.

\noindent a) $H=\Gm$, $U$ is a nontrivial form of $\Ga$ (cf. App. III, \S\ 5).

Let $k$ be an imperfect field of characteristic $2$, $u$ an element of $k$ such that $u^{1/2}$ does not belong to $k$. Consider the affine group $E$ with ring $k[X,Y,(X^2+uY^2)^{-1}]$, where multiplication is given by the comorphism
\[
(X,Y)\longmapsto(XX'+uYY',XY'+YX').
\]
The group $E$ is smooth, connected, commutative, of dimension $2$; the subscheme $Y=0$ defines a subgroup $H\simeq\Gm$. The kernel $K$ of raising to the second power in $E$ has equation
\[
X^2+uY^2=1,
\]
hence has dimension $1$. The group $K$ contains the unipotent radical of $E$ (defined over $\bar k$), but also the contribution of $H$, which is isomorphic to $\bmu_2$. Since $K$ is reduced over $k$, the unipotent radical of $E$ is not defined over $k$, and $U=E/H$ does not lift in $E$. (One immediately checks that $U$ is the form of $\Ga$ with ring $k[V,W]/(V+uV^2+W^2)$, the morphism $E\to U$ corresponding to the comorphism $V\mapsto Y^2/(X^2+Y^2u)$, $W\mapsto XY/(X^2+Y^2u)$.)

\noindent b) $H=\Gm$, $U=\mathbf Z/2\mathbf Z$, $k$ imperfect of characteristic $2$.

Choosing $k$ and $u$ as in a), consider the subgroup $E$ of $\mathrm{GL}_2$ generated by the element $X$ such that
\[
X=\begin{pmatrix}0&1\\u&0\end{pmatrix}
=\begin{pmatrix}u^{1/2}&0\\0&u^{1/2}\end{pmatrix}
\begin{pmatrix}0&u^{-1/2}\\u^{1/2}&0\end{pmatrix}.
\]
The group $E$ is an extension of $\mathbf Z/2\mathbf Z$ by $\Gm$, but this extension is not trivial, since the unipotent part of $X$ is not defined over $k$.

\noindent c) $H=\bmu_p$, $U=\balpha_p$, $k$ imperfect.

Let $\mathfrak e$ be the commutative Lie $p$-algebra generated by two elements $X$ and $Y$ such that $X^{(p)}=X$ and $Y^{(p)}=aX$. By App. II 2.2, $\mathfrak e$ is the Lie $p$-algebra of an algebraic group $E$, an extension of $\balpha_p$ by $\bmu_p$, but this extension is trivial if and only if there exists $b\in k$ such that $b^p=a$ (for one then has $(bX+Y)^{(p)}=0$).

\noindent d) $H=\bmu_2$, $U=\balpha_2\times\balpha_2$, $E$ noncommutative, $k$ a field of characteristic $2$.

Consider the special linear group $\mathrm{SL}_{2,k}$ and let $E={}_{F}(\mathrm{SL}_{2,k})$. The group $E$ is a radicial group of height $1$, whose Lie algebra is generated by three elements $X,Y,Z$ satisfying the following relations:
\[
[X,Y]=Z,\qquad[X,Z]=[Y,Z]=0,
\]
\[
X^{(p)}=Y^{(p)}=0,\qquad Z^{(p)}=Z.
\]
Consequently, $E$ is a central extension of $U\simeq\balpha_2\times\balpha_2$ by $\bmu_2$. Each factor $\balpha_2$ of $U$ lifts uniquely in $E$, but $U$ itself does not lift in $E$, since $[X,Y]\ne0$.

% Source PDF p. 39.
\sgasection{7}{Nilpotent affine algebraic groups}
\numpara{7.1}\label{XVII.7.1}\textbf{Extensions of groups of multiplicative type.}
\begin{proposition}{7.1.1}\label{XVII.7.1.1}
Let $S$ be a prescheme, $H$ and $K$ two group $S$-preschemes of multiplicative type and of finite type, $E$ a group prescheme which is an extension of $K$ by $H$ (i.e. $K$ is the quotient of $E$ by $H$ for the fpqc topology). Then $E$ is of multiplicative type in the following two cases:

\noindent a) $E$ is commutative.

\noindent b) The fibers of $K$ are connected.
\end{proposition}
\par\noindent\emph{Proof. i) Case where $S$ is the spectrum of a field $k$.}
The assertion to be proved is local for the fpqc topology, which allows us to suppose $k$ algebraically closed, hence $H$ and $K$ diagonalizable. Note that $K$ acts trivially on $H$ by inner automorphisms: this is clear in case a), and follows from 6.2.1 in case b). By induction on the length of a suitable composition series of $K$, one reduces to the case where $K$ is of one of the following three types: a) $K=\Gm$, b) $K=\bmu_p$, c) $K=\bmu_q$ with $(q,p)=1$ and, in this case, $E$ is commutative. Using now an embedding of $H$ in $(\Gm)^r$, one deduces that $E$ embeds in an extension of $K$ by $(\Gm)^r$. One may therefore suppose $H=(\Gm)^r$.

\noindent a) If $K=\Gm$, $E$ is a smooth connected affine algebraic group (Exp. VIB 9.2) of unipotent rank zero, hence is a torus.

\noindent b) $K=\bmu_p$. In this case, $E$ is a trivial extension. Indeed, since $H$ is smooth, the canonical morphism $\Lie E\to\Lie\bmu_p$ is surjective (App. II 3.2), and it suffices to apply 5.8.2 i), taking account of App. II 2.2.

\noindent c) $K=\bmu_q$, with $(q,p)=1$ and $E$ commutative. Here again, the extension $E$ is trivial. Indeed, let $x\in E(k)$ be a lifting of a generator $\bar x$ of $\bmu_q(k)$. The element $x^q$ is an element of $H(k)$, hence is of the form $y^q$, $y\in H(k)$ (note that $(\Gm)^r(k)$ is $q$-divisible). Since $E$ is commutative, $y^{-1}x$ is a lifting of $\bar x$ of order $q$.

\par\noindent\emph{ii) General case.}
The groups $H$ and $K$ are flat, affine and of finite presentation over $S$ (Exp. IX 2.1), and consequently the same holds for $E$ (Exp. VIB 9.2). Using then the general technique of VIB \S\ 10, we reduce to the case where $S$ is noetherian. To show that $E$ is of multiplicative type, it then suffices to prove that $E$ is commutative and that ${}_nE$ is finite over $S$ for every $n$ (Exp. X 4.8 b)).
% typo: missing closing parenthesis in the reference to X 4.8 b); supplied.

\noindent a) $E$ is commutative. One must verify that the morphism
\[
E\times_S E\longrightarrow E,\qquad(x,y)\longmapsto[x,y]=xyx^{-1}y^{-1}
\]
factors through the unit section of $E$, and it suffices to verify this when $S$ is the spectrum of an artinian local ring. But then $E$ is of multiplicative type by i) and Exp. X 2.3.

% Source PDF p. 40.
\noindent b) ${}_nE$ is finite over $S$. Indeed, one has the exact sequence
\[
0\longrightarrow{}_nH\longrightarrow{}_nE\longrightarrow{}_nK\xrightarrow{u}H_n
\]
(where $H_n$ is the cokernel of raising to the $n$th power in $H$). One knows that $H_n$ is of multiplicative type (Exp. IX 2.7), hence separated, and that ${}_nH$ and ${}_nK$ are finite over $S$; $\Ker u$ is a closed subgroup of ${}_nK$, hence is finite over $S$, and ${}_nE$ is finite over $S$ as an extension of a finite group by a finite group (Exp. VIB 9.2).

\numpara{7.2}\label{XVII.7.2}\textbf{Structure of commutative affine algebraic groups.}
\begin{theorem}{7.2.1}\label{XVII.7.2.1}
Let $k$ be a field, $G$ a commutative affine algebraic $k$-group. Then:

\noindent a) $G$ contains a largest subgroup of multiplicative type $M$. The group $M$ is characteristic in $G$ and $G/M$ is unipotent, and its formation commutes with extension of the field $k$.

\noindent b) If $k$ is perfect, $G$ is the direct product of $M$ and a unipotent algebraic subgroup $U$, and this uniquely.
\end{theorem}
\par\noindent\emph{Proof. i) $k$ algebraically closed.}
When $G$ is smooth, 7.2.1 b) is well known (BIBLE \S\ 4 Th. 4). If $G$ is radicial of height $1$, the decomposition of $\Lie G$ described in 4.2.2 corresponds, taking account of 4.3.1 v) and App. II 2.2, to a decomposition of $G$ of the type 7.2.1 b). In the general case, $G$ admits a composition series whose successive quotients are smooth or radicial of height $1$ (App. II 3.1). To prove 7.2.1 b), it then suffices to note that if one has an exact sequence of commutative algebraic groups
\[
0\longrightarrow G'\longrightarrow G\longrightarrow G''\longrightarrow0,
\]
where $G'$ (resp. $G''$) is the product of a group of multiplicative type and a unipotent group $G'=M'\cdot U'$ (resp. $G''=M''\cdot U''$), then the same holds for $G$. Indeed, consider the exact sequence
\[
0\longrightarrow G'/U'\longrightarrow G/U'\longrightarrow G''\longrightarrow0.
\]
By 7.1.1 a), the inverse image in $G/U'$ of $M''$ is a subgroup of multiplicative type $M_1$. The group $M_1$ lifts to a subgroup $M$ of $G$ (5.1.1 i) a)). Similarly, using this time 6.1.1 C), one proves that there exists a unipotent subgroup $U$ of $G$ which is an extension of $U''$ by $U'$. It is clear that $G=M\cdot U$.

\par\noindent\emph{ii) $k$ arbitrary.}
By i), $G_{\bar k}$ is the direct product of a group of multiplicative type $M_{\bar k}$ with a unipotent group $U_{\bar k}$. Put $S=\bar k\times_k\bar k$, and let $M_1$ and $M_2$ be the two inverse images of $M_{\bar k}$ by the two projections $S\rightrightarrows\bar k$. The group $G_S/M_2=(G/M_{\bar k})_S$ has unipotent fibers, so the image of $M_1$ in $G_S/M_2$ is zero (2.4 i)), and $M_2$ contains $M_1$. Similarly, $M_1$ contains $M_2$, and finally $M_1=M_2$. By fpqc descent, it follows that $M_{\bar k}$ comes from an algebraic subgroup $M$ of $G$. It is clear that $M$ is of multiplicative type, that $G/M$ is unipotent, and that formation of $M$ is compatible with every extension of the field $k$. For every $k$-prescheme $S$, every subgroup of multiplicative type $H$ of $G_S$ is contained in $M_S$. Indeed, by 2.5 its image in the group with unipotent fibers $(G/M)_S=G_S/M_S$ is zero. Taking in particular for $H$ the transform of $M_S$ by an automorphism of $G_S$, one deduces that $M$ is characteristic in $G$.

% Source PDF p. 41.
Finally, if $k$ is perfect, $G/M$ lifts in $G$ to a unipotent group, and this uniquely by 6.1.1 C).
\begin{remark}{7.2.2}\label{XVII.7.2.2}
\noindent i) If $k$ is imperfect, the unipotent component of $G_{\bar k}$ is not necessarily defined over $k$, as shown by example 6.4 a).

\noindent ii) Contrary to what happens for the component of multiplicative type $M$, the unipotent component $U$ is in general not characteristic in $G$ (and this regardless of the characteristic of $k$). Of course, uniqueness of the decomposition 7.2.1 b) implies that $U$ is invariant under every $k$-automorphism of $G$. But if $U$ and $M$ are such that there exist a $k$-prescheme $S$ and a nonzero $S$-homomorphism $h:U_S\to M_S$ (cf. 2.6), one deduces an $S$-automorphism of $G_S$, $(u,m)\mapsto(u,h(u)+m)$, which does not leave $U_S$ invariant.

\noindent iii) If $G$ is finite over $k$, $G/M$ corresponds by Cartier duality (2.6) to the connected component of the dual $D(G)$ of $G$.
\end{remark}
\numpara{7.3}\label{XVII.7.3}\textbf{Structure of nilpotent affine algebraic groups.}
\begin{theorem}{7.3.1}\label{XVII.7.3.1}
Let $k$ be a field, $G$ a nilpotent algebraic $k$-group (Exp. VIB \S\ 8), affine, connected. Then $G$ has a largest subgroup of multiplicative type $M$. The group $M$ is central and characteristic, and $G/M$ is a unipotent algebraic group.
\end{theorem}
Let $Z$ be the center of $G$, $M$ the largest subgroup of multiplicative type of $Z$ (7.2.1). Since $Z$ is characteristic in $G$ and $M$ characteristic in $Z$, $M$ is characteristic in $G$. It suffices to show that $G/M$ is unipotent. By induction on the length of the ascending central series of $G$, this will follow more generally from the following lemma:
\begin{lemma}{7.3.2}\label{XVII.7.3.2}
(Rosenlicht). Let $G$ be a connected algebraic $k$-group, $Z$ its center. Then the center $Z'$ of $G'=G/Z$ is unipotent.
\end{lemma}
\par\noindent\emph{Proof.}
We may suppose $k$ algebraically closed. It then suffices to show that $Z'$ contains no subgroup isomorphic to $\bmu_\ell$ for any prime number $\ell$ (4.6.1 vi)). Let therefore $\bmu_\ell$ be a subgroup of $Z'$, $N$ its inverse image in $G$. Since $\bmu_\ell$ is central in $G'$, $N$ is invariant in $G$.

\par\noindent\emph{i) Case where $(\ell,p)=1$.}
One can find an element $x$ of $N(k)$ and an integer $n$ having the following properties:

\noindent a) $x$ lifts a generator $\bar x$ of $\bmu_\ell$.

\noindent b) $x^{\ell^n}\in Z^0(k)$;

(it suffices to choose a lifting of $\bar x$ whose image in $N/Z^0(k)$ belongs to the Sylow $\ell$-subgroup).

Raising to the $\ell$th power in the commutative group $Z^0$ is an étale morphism, so $Z^0(k)$ is $\ell$-divisible. Consequently, after multiplying $x$ by an element of $Z^0(k)$, one may suppose that $x^{\ell^n}=0$. The group $N$ is then generated by two commutative groups which commute ($Z$ and the group generated by $x$), hence is commutative. The group ${}_{\ell^n}N$ is a group of multiplicative type, characteristic in $N$, hence invariant in $G$, and consequently central, $G$ being connected (Exp. IX 5.5). Thus ${}_{\ell^n}N$ is contained in $Z$, which contradicts the fact that its image in $G'$ contains $\bmu_\ell$.
% typo?: x^(ell^n)=0 in multiplicative notation is kept as printed.

% Source PDF p. 42.
\noindent ii) $\ell=p$. There then exists an integer $n$ such that the image of ${}_{F^n}N$ in $G'$ contains $\bmu_p$, so that one has the exact sequence
\[
1\longrightarrow K\longrightarrow {}_{F^n}N\longrightarrow\bmu_p\longrightarrow1.\tag{*}
\]
The group $K$ is contained in $Z$, hence is commutative; it then follows from 7.2.1 and 7.1.1 b) and 5.5.1 that there exists a subgroup of multiplicative type contained in ${}_{F^n}N$ whose image in the quotient by $K$ is $\bmu_p$. One deduces as in i) that ${}_{F^n}N$ is commutative. The component of multiplicative type of ${}_{F^n}N$ (7.2.1) is characteristic in ${}_{F^n}N$, hence invariant in $G$, hence central (Exp. IX 5.5); since its image in $G'$ contains $\bmu_p$, one obtains a contradiction.

\sgasection{A}{Appendix I. Hochschild cohomology and extensions of algebraic groups}
\numpara{A.1}\label{XVII.A.1}\textbf{Definition of the cohomology groups.}
Let $k$ be a field, $G$ an algebraic $k$-group, $\Abcat_G$ the abelian category of commutative algebraic $k$-groups on which $G$ acts. If $A\in\Ob\Abcat_G$, the functor $h_A:(\Sch/k)^\circ\to\Ens$ canonically defined by $A$ is a $G$-$\mathbf Z$-module in the sense of I 3.2. We may therefore consider the standard complex $C^\bullet(G,A)$ of algebraic cochains of $G$ with values in $A$ (Exp. I 5.1), as well as the group of $i$-cocycles $Z^i(G,A)$, of $i$-coboundaries $B^i(G,A)$ and the $i$th cohomology group $H^i(G,A)$. As usual, $H^0(G,A)$ is identified with the group $A^G(k)$, where $A^G$ is the $k$-functor of invariants of $A$ under $G$. The group $H^1(G,A)$ classifies trivial principal homogeneous spaces under $A$ on which $G$ acts.

The functor $A\mapsto H^\bullet(G,A)$ is in general not a cohomological functor from the category $\Abcat_G$ with values in $\Abcat$; however, one has the following proposition:
\begin{proposition}{A.1.1}\label{XVII.A.1.1}
Let $0\to A'\xrightarrow{u}A\xrightarrow{v}A''\to0$ be an exact sequence in $\Abcat_G$. Then:

\noindent a) If $v$ has a section (that is to say, if there exists a $k$-morphism of preschemes $s:A''\to A$ such that $vs=1_{A''}$), one has the usual exact cohomology sequence
\[
\cdots\longrightarrow H^i(G,A)\longrightarrow H^i(G,A'')\xrightarrow{d}H^{i+1}(G,A')\longrightarrow\cdots
\]
\noindent b) If $A(k)\to A''(k)$ is surjective, one has the exact sequence
\[
0\to A'^G(k)\to A^G(k)\to A''^G(k)\xrightarrow{d}H^1(G,A')\to H^1(G,A)\to H^1(G,A'').\tag{1}
\]
\end{proposition}
\par\noindent\emph{Proof.}
a) One notes that existence of a section implies exactness of the sequence of complexes
\[
0\longrightarrow C^\bullet(G,A')\longrightarrow C^\bullet(G,A)\longrightarrow C^\bullet(G,A'')\longrightarrow0.
\]
% Source PDF p. 43.
b) If $x''\in A''^G(k)$, its inverse image in $A$ is a trivial principal homogeneous space under $A'$ (since $A(k)\to A''(k)$ is supposed surjective) on which $G$ acts, hence defines an element $d(x'')\in H^1(G,A')$. Exactness of sequence (1) is then immediate.

\numpara{A.2}\label{XVII.A.2}\textbf{The group $\Ext_{\mathrm{alg}}(G,A)$.}
Let $A$ and $G$ be two algebraic $k$-groups, $E$ and $E'$ two algebraic $k$-groups which are extensions of $G$ by $A$. These two extensions are said to be isomorphic if there exists a group $k$-morphism $u:E\to E'$ making the diagram commutative:
\[
\xymatrix{&E\ar[dr]\ar[dd]^u&\\A\ar[ur]\ar[dr]&&G\\&E'\ar[ur]&}
\]
The group $E$ acts on $A$ by inner automorphisms, and if $A$ is commutative, this action factors through $G$, so $G$ acts on $A$. Conversely, if $A\in\Ob\Abcat_G$, we denote by $\Ext_{\mathrm{alg}}(G,A)$ the set of classes of algebraic extensions $E$ of $G$ by $A$ for which the action of $G$ on $A$ defined by $E$ and that coming from the structure of an object of $\Abcat_G$ coincide.

$\Ext_{\mathrm{alg}}(G,A)$ is naturally a bifunctor covariant in $A$ and contravariant in $G$. More precisely:

\noindent a) If $f:A\to B$ is a morphism in $\Abcat_G$ and $E$ represents an element of $\Ext_{\mathrm{alg}}(G,A)$, one defines $f_*(E)\in\Ext_{\mathrm{alg}}(G,B)$ as the class of the extension of $G$ by $B$ equal to the quotient of the semidirect product $B\cdot E$ ($E$ acting on $B$ through $G$) by the algebraic subgroup which is the image of $A$ under the morphism $(f,i)$ (this quotient is representable by Exp. VIA \S\ 5), so that one has a commutative diagram
\[
\xymatrix{1\ar[r]&A\ar[r]^i\ar[d]_f&E\ar[r]\ar[d]&G\ar[r]\ar[d]^{1_G}&1\\1\ar[r]&B\ar[r]&f_*(E)\ar[r]&G\ar[r]&1.}
\]
\noindent b) If $g:H\to G$ is a $k$-morphism of algebraic $k$-groups and $E$ is an extension of $G$ by $A$, the fibered product $E\times_G H$ is naturally an extension of $H$ by $A$, denoted by $g^*(E)$. One therefore has a commutative diagram
\[
\xymatrix{1\ar[r]&A\ar[r]\ar[d]_{1_A}&g^*(E)\ar[r]\ar[d]&H\ar[r]\ar[d]^g&1\\1\ar[r]&A\ar[r]&E\ar[r]&G\ar[r]&1.}
\]
Adapting the proofs given in J.-P. Serre, \emph{Groupes algébriques et corps de classe}, chap. VII, one equips $\Ext_{\mathrm{alg}}(G,A)$ with a natural abelian group structure, functorial in $A$ and $G$.
% typo?: the pushout construction prints (f,i), without a minus sign; kept.

% Source PDF p. 44.
\begin{proposition}{A.2.1}\label{XVII.A.2.1}
Let $0\to A'\to A\to A''\to0$ be an exact sequence in $\Abcat_G$; then:

\noindent a) One has a canonical exact sequence of abelian groups
\[
Z^1(G,A)\to Z^1(G,A'')\xrightarrow{d}\Ext_{\mathrm{alg}}(G,A')\to\Ext_{\mathrm{alg}}(G,A)\to\Ext_{\mathrm{alg}}(G,A'').
\]
\noindent b) If $A(k)\to A''(k)$ is surjective, one deduces from a) the exact sequence
\[
H^1(G,A)\to H^1(G,A'')\xrightarrow{d}\Ext_{\mathrm{alg}}(G,A')\to\Ext_{\mathrm{alg}}(G,A)\to\Ext_{\mathrm{alg}}(G,A'').\tag{2}
\]
\end{proposition}
The exact sequence of a) generalizes the usual exact sequence of the $\Hom(\cdot,\cdot)$ valid in the context of commutative extensions (loc. cit.) and is proved in the same way. Let us simply recall the definition of the coboundary $d:Z^1(G,A'')\to\Ext_{\mathrm{alg}}(G,A')$. For this, consider the extension
\[
1\longrightarrow A'\longrightarrow A\cdot G\longrightarrow A''\cdot G\longrightarrow1,
\]
deduced in an evident way from the exact sequence $0\to A'\to A\to A''\to0$. If $u\in Z^1(G,A'')$, $u$ defines in the usual way a section homomorphism $u:G\to A''\cdot G$. One then has $d(u)=u^*(A\cdot G)$.

\numpara{A.3}\label{XVII.A.3}\textbf{Comparison of $H^2(G,A)$ and $\Ext_{\mathrm{alg}}(G,A)$.}
It is well known, in the case of abstract groups, that there exists a functorial isomorphism between the abelian groups $H^2(G,A)$ and $\Ext(G,A)$. Similarly, in the present case, if $A$ is an element of $\Abcat_G$, to every $2$-cocycle $u\in Z^2(G,A)$ one can associate an algebraic group structure on the prescheme $A\times_k G$ making it an element of $\Ext_{\mathrm{alg}}(G,A)$. Moreover, this extension is trivial if and only if $u\in B^2(G,A)$ (cf. Exp. III 1.2.2). Recall that the composition law on $A\times G$ is defined by the formula
\[
(a,g)(a',g')=a+{}^ga'+u(g,g').
\]
% typo?: the printed formula omits the second component gg'; kept.
It is clear that the extensions of $G$ by $A$ thus obtained are not arbitrary, since they have a section. But conversely, if $E\in\Ext_{\mathrm{alg}}(G,A)$ has a section $s$, $E$ is isomorphic to the extension of $G$ by $A$ associated to the $2$-cocycle $u$ such that
\[
u(g,g')=s(g)s(g')s(gg')^{-1}.
\]
One finally obtains the following proposition:
\begin{proposition}{A.3.1}\label{XVII.A.3.1}
There exists a functorial isomorphism between the bifunctors with values in abelian groups
\[
(G,A)\longmapsto H^2(G,A)\qquad\text{and}\qquad(G,A)\longmapsto\Ext_s(G,A),
\]
where $\Ext_s(G,A)$ denotes the subgroup of $\Ext_{\mathrm{alg}}(G,A)$ formed by the classes of extensions of $G$ by $A$ which have a section.
\end{proposition}

\sgasection{B}{Appendix II. Recollections and complements on radicial groups}
Let $p$ be a prime number $>1$, and let $S$ be an $\mathbf F_p$-prescheme.

% Source PDF p. 45.
\numpara{B.1}\label{XVII.B.1}\textbf{The Frobenius morphism.}
For every $S$-prescheme $X$ and every integer $n>0$, denote by $P^n:X\to X$ the $\mathbf F_p$-endomorphism corresponding to raising to the $p^n$th power in $\OO_X$, and denote by $X^{(n)}$ the $S$-prescheme which is the inverse image of $X$ under the morphism $P^n:S\to S$. There then exists a unique $S$-morphism $F^n:X\to X^{(n)}$ making the diagram commutative:
\[
\xymatrix@C=3em@R=3em{&&X\ar[dll]_{P^n}\ar[dl]^{F^n}\ar[ddl]\\X\ar[d]&X^{(n)}\ar[l]\ar[d]&\\S&S\ar[l]^{P^n}&}
\]
It is clear that $F^n$ is identified with the ``$n$th iterate'' of $F^1=F$, called the Frobenius morphism of $X/S$.

If $G$ is a group $S$-prescheme, $G^{(n)}$ is a group $S$-prescheme, and $F^n:G\to G^{(n)}$ is a group $S$-morphism. Its kernel ${}_{F^n}(G)$ is a characteristic group subprescheme of $G$ (i.e. stable under the functor $\Aut_{S\text{-gr}}(G)$), radicial over $S$. If $G$ is a radicial group $S$-prescheme, one says that $G$ has height $\leq h$ if ${}_{F^h}(G)=G$.

\numpara{B.2}\label{XVII.B.2}\textbf{Groups and Lie $p$-algebras.}
If $G$ is a group $S$-prescheme, $\Lie(G)$ (Exp. II) is naturally a restricted Lie $p$-Algebra (Exp. VIIA \S\S\ 5 and 6). In particular, one has the following result (cf. Exp. VIIA):
\setcounter{footnote}{10}
\begin{proposition}{B.2.1}\label{XVII.B.2.1}
\noindent i) $\Lie(\balpha_p)_S=\Lie(\Ga)_S$\nde{The definition of $\balpha_p$, already given at the beginning of this exposé, has been deleted here.} is a sheaf of $\OO_S$-modules, free over $\OO_S$ of rank $1$, generated by an element $X$ such that $X^{(p)}=0$.

\noindent ii) $\Lie(\bmu_p)_S=\Lie(\Gm)_S$ is a sheaf of $\OO_S$-modules, free over $\OO_S$ of rank $1$, with a canonical basis $X$ such that $X^{(p)}=X$.
\end{proposition}
Let us now recall the fundamental result proved in Exp. VIIA \S\ 7:
\begin{theorem}{B.2.2}\label{XVII.B.2.2}
Suppose $S$ affine with ring $A$. Then the functor
\[
G\longmapsto\Lie G
\]
establishes an equivalence of categories between, on the one hand, the category of group $S$-preschemes $G$ of finite presentation and flat over $S$, of height $1$, whose Lie algebra is locally free over $\OO_S$, and, on the other hand, the category of restricted Lie $p$-$A$-algebras locally free of finite rank.

% Source PDF p. 46.
Moreover, if $G$ is as above and $H$ is a group $S$-prescheme of finite presentation, the canonical morphism
\[
\Hom_{S\text{-gr}}(G,H)\longrightarrow\Hom_{p\text{-}A\text{-Lie}}(\Lie G(S),\Lie H(S))
\]
is an isomorphism.
\end{theorem}
\numpara{B.3}\label{XVII.B.3}\textbf{Radicial groups and smooth groups.}
We now suppose that $S$ is the spectrum of a field $k$ of characteristic $p$.

Recall that in Exp. VIA \S\ 5, it was shown that if $G$ is an algebraic $k$-group and $H$ an algebraic subgroup of $G$, then the sheaf $G/H$ (a sheaf for the fpqc topology) is representable. Let us then recall (VIIA 8.3):
\begin{proposition}{B.3.1}\label{XVII.B.3.1}
Let $G$ be an algebraic $k$-group. Then there exists an integer $m$ such that for every $n\geq m$, the algebraic group $G/{}_{F^n}(G)$ is smooth over $k$.
\end{proposition}
\begin{proposition}{B.3.2}\label{XVII.B.3.2}
Consider an exact sequence of algebraic $k$-groups
\[
1\longrightarrow G'\longrightarrow G\xrightarrow{u}G''\longrightarrow1
\]
and the following assertions:

\noindent i) The morphism $u$ is smooth.

\noindent ii) $G'$ is smooth over $k$.

\noindent iii) For every integer $n>0$, one has the exact sequence
\[
1\longrightarrow {}_{F^n}(G')\longrightarrow {}_{F^n}(G)\longrightarrow {}_{F^n}(G'')\longrightarrow1.
\]
\noindent iv) The morphism ${}_{F}G\to {}_{F}G''$ is an epimorphism.

\noindent v) The morphism $(\Lie G)(k)\to(\Lie G'')(k)$ is surjective.

Then one has the following implications:
\[
\text{i)}\Longleftrightarrow\text{ii)}\Longrightarrow\text{iii)}\Longrightarrow\text{iv)}\Longleftrightarrow\text{(v)}.
\]
Moreover, if $G$ is smooth, the five assertions are equivalent.
\end{proposition}
$\text{i)}\Leftrightarrow\text{ii)}$ by Exp. VIB 9.2 vii).

$\text{ii)}\Rightarrow\text{iii)}$. If $G'$ is smooth, $F^n:G'\to G'^{(n)}$ is an epimorphism, and iii) follows from the ``snake'' diagram.

$\text{iii)}\Rightarrow\text{iv)}$ is clear.

$\text{iv)}\Leftrightarrow\text{v)}$ by Theorem B.2.2.

$\text{v)}\Rightarrow\text{ii)}$ when $G$ is smooth. Indeed, $G''$ is then smooth, and v) implies that one has $\dim G'=\dim_k(\Lie G')(k)$, so $G'$ is smooth.

\sgasection{C}{Appendix III. Remarks and complements concerning exposés XV, XVI, XVII}
\numpara{C.1}\label{XVII.C.1}
It may be that Propositions 1.2 and 1.2 bis of XV remain true if one removes the hypothesis that $H_0$ is smooth over $S_0$. This is in particular the case if $G$ is finite, flat and commutative.
